\documentclass[12pt,reqno]{amsart}

\usepackage{amsfonts,amsmath,amssymb,amsthm,mathtools}
\usepackage{mathrsfs}       % \mathscr
\usepackage{xfrac}
\usepackage{enumerate}
\usepackage{xcolor}
\usepackage[hidelinks]{hyperref}
\usepackage{microtype}

\linespread{1.1}
\allowdisplaybreaks
\numberwithin{equation}{subsection}

% ---- Input all macros ---------------------------------------------------
% All notation lives here. Add symbols with a brief comment; never define
% notation inline in a section.

% ---- Von Neumann algebras / operators (starter set) ---------------------
\newcommand{\M}{\mathcal{M}}
\newcommand{\N}{\mathcal{N}}
\newcommand{\B}{\mathcal{B}}
\def\H{\mathcal{H}}                  % \H is otherwise a text accent; override it
\def\K{\mathcal{K}}
\newcommand{\comm}[1]{#1'}          % commutant
\newcommand{\bicomm}[1]{#1''}       % double commutant

% ---- Number fields ------------------------------------------------------
\newcommand{\Cplx}{\mathbb{C}}
\newcommand{\Real}{\mathbb{R}}
\newcommand{\Intg}{\mathbb{Z}}
\newcommand{\Nat}{\mathbb{N}}
\newcommand{\Rat}{\mathbb{Q}}

% ---- Polynomial orders -------------------------------------------------
\newcommand{\polyord}{\operatorname{ord}}

% ---- Common shorthands --------------------------------------------------
\newcommand{\ip}[2]{\langle #1, #2 \rangle}
\newcommand{\norm}[1]{\left\| #1 \right\|}
\newcommand{\eps}{\varepsilon}


% ---- Theorem environments -----------------------------------------------
% Section-numbered family (single counter: thm)
\theoremstyle{plain}
\newtheorem{thm}{Theorem}[section]
\newtheorem{lem}[thm]{Lemma}
\newtheorem{cor}[thm]{Corollary}
\newtheorem{prop}[thm]{Proposition}
\newtheorem{claim}[thm]{Claim}

\theoremstyle{definition}
\newtheorem{definition}[thm]{Definition}
\newtheorem{example}[thm]{Example}
\newtheorem{rem}[thm]{Remark}
\newtheorem{nota}[thm]{Notation}
\newtheorem{ques}[thm]{Question}

% Letter-numbered family (for main results stated in the introduction)
\newtheorem{letterthm}{Theorem}
\renewcommand{\theletterthm}{\Alph{letterthm}}
\newtheorem{lettercor}[letterthm]{Corollary}
\renewcommand{\thelettercor}{\Alph{lettercor}}
\newtheorem{letterconj}[letterthm]{Conjecture}
\renewcommand{\theletterconj}{\Alph{letterconj}}

% Unnumbered environments
\newtheorem*{mainthm}{Main Theorem}
\newtheorem*{question}{Question}
\newtheorem*{opr}{Open Problem}

% Step environment (inside long proofs)
\newtheorem{step}{Step}

% ---- Front matter -------------------------------------------------------
\title[Periodic Blocks in Rational Riordan Arrays]
{Periodic Blocks in Rational Riordan Arrays over Finite Fields}

\author{Nikolai Krylov}
\address{}
\email{}

\keywords{Riordan array, finite field, periodic sequence, circulant matrix}
\subjclass[2020]{05A15, 11B85, 12E20}

\begin{document}

\begin{abstract}
We study Riordan arrays over a finite field $\mathbb F_q$ whose defining pair
consists of rational functions. Their columns have eventually periodic
coefficient sequences. We compute the least eventual periods under a
coprimality hypothesis and show that consecutive periods are equal or differ
by a factor of the characteristic. Periodic blocks of neighboring columns
satisfy a linear equation in a circulant quotient ring. The coefficient
multiplying the later block is generally singular, so this equation defines a
correspondence rather than an invertible evolution. We first give the theory
over $\mathbb F_5$, following a Lean formalization, and then extend the
statements to arbitrary finite fields. Pascal's array provides the basic
example.
\end{abstract}

\maketitle

\section{Introduction}

% State the setting and the main results. Promote settled material here from
% dialogue/dialogue.md.

\section{The formalized case over \texorpdfstring{$\mathbb F_5$}{F5}}

This section follows the Lean formalization. All polynomials and power series
have coefficients in $\mathbb F_5=\mathbb Z/5\mathbb Z$. The notation is
chosen to agree with the formal definitions.

\subsection{Polynomial orders and column periods}

For $Q\in\mathbb F_5[t]$, define
\[
\polyord(Q)=\polyord_{\mathbb F_5[t]/(Q)}(t),
\]
the multiplicative order of the residue class of $t$. When $Q$ is nonzero and
$Q(0)\ne0$, this is the least positive integer $N$ such that
\begin{equation}\label{eq:f5-order}
Q(t)\mid t^N-1.
\end{equation}
More generally, the formalization proves the divisibility criterion
\begin{equation}\label{eq:f5-order-divisibility}
\polyord(Q)\mid N
\quad\Longleftrightarrow\quad
Q(t)\mid t^N-1.
\end{equation}

Fix
\[
g(t)=\frac{p_1(t)}{p_2(t)},
\qquad
f(t)=\frac{p_3(t)}{p_4(t)},
\]
where $p_2(0)p_4(0)\ne0$. The generating function of column $k$ is
\begin{equation}\label{eq:f5-column}
H_k(t)=\frac{p_1(t)p_3(t)^k}{p_2(t)p_4(t)^k}.
\end{equation}
For an ordinary Riordan array one additionally assumes $g(0)\ne0$,
$f(0)=0$, and $f'(0)\ne0$. The formalized periodicity statements themselves
use only the displayed rational column generating functions and the
nonvanishing denominator constants.
Set
\[
Q_k(t)=p_2(t)p_4(t)^k,
\qquad
\pi_k=\polyord(Q_k).
\]

\begin{prop}\label{prop:f5-eventual}
The coefficient sequence of $H_k$ is eventually $N$-periodic whenever
$\pi_k\mid N$. If
\[
\gcd\bigl(Q_k,p_1p_3^k\bigr)=1,
\]
then every eventual period of $H_k$ is divisible by $\pi_k$. In particular,
$\pi_k$ is its least eventual period.
\end{prop}

The first statement follows by multiplying Equation~\eqref{eq:f5-column} by
$t^N-1$. For the converse, eventual $N$-periodicity makes
$(t^N-1)H_k$ a polynomial. Coprimality then implies
$Q_k\mid t^N-1$, and Equation~\eqref{eq:f5-order-divisibility} gives
$\pi_k\mid N$.

\subsection{The five-adic period tower}

The formalized argument uses
\begin{equation}\label{eq:f5-frobenius}
(t^N-1)^5=t^{5N}-1.
\end{equation}
Since $Q_k\mid Q_{k+1}$, one has $\pi_k\mid\pi_{k+1}$. If $k\geqslant1$,
then
\[
Q_{k+1}=p_2p_4^{k+1}\mid(p_2p_4^k)^2=Q_k^2.
\]
It follows from Equation~\eqref{eq:f5-frobenius} that
$Q_{k+1}\mid t^{5\pi_k}-1$, and hence $\pi_{k+1}\mid5\pi_k$.

\begin{prop}\label{prop:f5-ratio}
Suppose $p_2$ and $p_4$ are nonzero with nonzero constant terms. If
$k\geqslant1$, then
\[
\pi_{k+1}=\pi_k
\qquad\text{or}\qquad
\pi_{k+1}=5\pi_k.
\]
Consequently, there is a nondecreasing function $u\colon\mathbb N\to\mathbb
N$ with $u(1)=0$ such that
\[
\pi_k=\pi_1 5^{u(k)}
\qquad(k\geqslant1).
\]
\end{prop}

The formalization also verifies the factorization formula. Define
\[
\lambda_5(m)=\min\{j\geqslant0:m\leqslant5^j\}
\qquad(m\geqslant1).
\]
If $\varphi$ is irreducible and $\varphi(0)\ne0$, then
\begin{equation}\label{eq:f5-prime-power}
\polyord(\varphi^m)=\polyord(\varphi)5^{\lambda_5(m)}.
\end{equation}
Suppose a common finite family of pairwise nonassociated irreducibles gives
\[
p_2=\prod_{i\in I}\varphi_i^{a_i},
\qquad
p_4=\prod_{i\in I}\varphi_i^{b_i},
\]
with $\varphi_i(0)\ne0$ and $a_i+kb_i\geqslant1$. Then
\begin{align}
\pi_k
&=\mathop{\operatorname{lcm}}_{i\in I}
  \left(\polyord(\varphi_i)5^{\lambda_5(a_i+kb_i)}\right)\label{eq:f5-lcm}\\
&=E5^{s_k},\label{eq:f5-explicit}
\end{align}
where
\[
E=\mathop{\operatorname{lcm}}_{i\in I}\polyord(\varphi_i),
\qquad
s_k=\max_{i\in I}\lambda_5(a_i+kb_i).
\]
The second equality uses the fact that each $\polyord(\varphi_i)$ is relatively
prime to five.

\subsection{Circulant blocks}

Let a power series $H(t)=\sum_{n\geqslant0}c_nt^n$ be eventually
$N$-periodic. For a starting index $M$ after the periodic tail has begun,
define
\[
B_{N,M}(H)=\sum_{r=0}^{N-1}c_{M+r}t^r
\quad\text{in}\quad
\mathcal C_N=\mathbb F_5[t]/(t^N-1).
\]
The formalization proves that taking a sufficiently late block intertwines
multiplication by a polynomial with multiplication by its class in
$\mathcal C_N$. It also proves that moving the starting index corresponds to
cyclic shift.

The column generating functions satisfy
\begin{equation}\label{eq:f5-neighbors}
p_4(t)H_{k+1}(t)=p_3(t)H_k(t).
\end{equation}

\begin{prop}\label{prop:f5-block}
Let $N=\pi_{k+1}$. There is an index $M_0$ such that, for every
$M\geqslant M_0$, the length-$N$ blocks read from columns $k$ and $k+1$
satisfy
\[
p_4(t)B_{N,M}(H_{k+1})=p_3(t)B_{N,M}(H_k)
\quad\text{in }\mathcal C_N.
\]
Moreover, $B_{N,M}(H_k)$ is the length-$\pi_k$ block repeated
$N/\pi_k$ times.
\end{prop}

If $p_4$ is nonconstant, its class in $\mathcal C_N$ is not a unit because
$p_4\mid Q_{k+1}\mid t^N-1$. Thus the displayed relation cannot generally be
solved by inverting $p_4$.

\subsection{Pascal's array modulo five}

For
\[
g(t)=\frac{1}{1-t},
\qquad
f(t)=\frac{t}{1-t},
\]
one has
\[
H_k(t)=\frac{t^k}{(1-t)^{k+1}},
\qquad
\pi_k=5^{\lambda_5(k+1)}.
\]
The first periods are
\[
1,5,5,5,5,25,\ldots,25,125,\ldots.
\]
For $N=\pi_{k+1}$, sufficiently late neighboring blocks satisfy
\[
(1-t)B_{k+1}=tB_k
\quad\text{in }\mathbb F_5[t]/(t^N-1).
\]
This is a cyclic discrete-antidifference equation. The formalization also
checks directly that the class of $1-t$ is not a unit in the corresponding
circulant ring.

\section{The extension to arbitrary finite fields}

Let $\mathbb F_q$ have characteristic $p$. The arguments of Section~2 extend
after replacing five by $p$. We state the results separately because the Lean
formalization currently treats only $\mathbb F_5$.

Fix nonzero polynomials $p_1,p_2,p_3,p_4\in\mathbb F_q[t]$ such that
$p_2(0)p_4(0)\ne0$, and impose $g(0)\ne0$, $f(0)=0$, and $f'(0)\ne0$. Set
\[
g(t)=\frac{p_1(t)}{p_2(t)},
\qquad
f(t)=\frac{p_3(t)}{p_4(t)}.
\]
The $k$th column generating function and its denominator before cancellation
are
\[
H_k(t)=\frac{p_1(t)p_3(t)^k}{p_2(t)p_4(t)^k},
\qquad
Q_k(t)=p_2(t)p_4(t)^k.
\]
Define
\[
\pi_k=\polyord(Q_k)
=\min\{N\geqslant1:Q_k(t)\mid t^N-1\}.
\]
The coefficients of $H_k$ are eventually $\pi_k$-periodic. If
$Q_k$ is coprime to $p_1p_3^k$, then $\pi_k$ is the least eventual period.

\subsection{Period growth}

The Frobenius identity is
\begin{equation}\label{eq:q-frobenius}
(t^N-1)^p=t^{pN}-1.
\end{equation}
As before, $Q_k\mid Q_{k+1}$ gives $\pi_k\mid\pi_{k+1}$. For
$k\geqslant1$, one also has
\[
Q_{k+1}\mid Q_k^2\mid Q_k^p.
\]
Together with Equation~\eqref{eq:q-frobenius}, this gives
$\pi_{k+1}\mid p\pi_k$.

\begin{prop}\label{prop:q-ratio}
If $k\geqslant1$, then
\[
\pi_{k+1}=\pi_k
\qquad\text{or}\qquad
\pi_{k+1}=p\pi_k.
\]
Consequently, there is a nondecreasing sequence $(u_k)_{k\geqslant1}$ with
$u_1=0$ such that
\[
\pi_k=\pi_1p^{u_k}.
\]
\end{prop}

For the explicit formula, define
\[
\lambda_p(m)=\min\{j\geqslant0:m\leqslant p^j\}
\qquad(m\geqslant1).
\]
Suppose a common finite family of pairwise nonassociated irreducibles gives
\[
p_2=\prod_{i\in I}\varphi_i^{a_i},
\qquad
p_4=\prod_{i\in I}\varphi_i^{b_i},
\]
where $\varphi_i(0)\ne0$ and $a_i+kb_i\geqslant1$. Write
$e_i=\polyord(\varphi_i)$. The roots of $\varphi_i$ lie in the multiplicative
group of a finite extension of $\mathbb F_q$, so $p\nmid e_i$. The
prime-power formula becomes
\[
\polyord(\varphi_i^m)=e_i p^{\lambda_p(m)}.
\]
Therefore
\begin{prop}\label{prop:q-formula}
Under the preceding factorization hypotheses,
\begin{align*}
\pi_k
&=\mathop{\operatorname{lcm}}_{i\in I}
  \left(e_i p^{\lambda_p(a_i+kb_i)}\right)\\
&=E p^{s_k},
\end{align*}
where
\[
E=\mathop{\operatorname{lcm}}_{i\in I}e_i,
\qquad
s_k=\max_{i\in I}\lambda_p(a_i+kb_i).
\]
\end{prop}

Thus rational second components produce a tower of changing block lengths.
This differs from the polynomial second-component setting, where the
denominator of the column generating function is fixed and one circulant
space suffices for all columns.

\subsection{The circulant correspondence}

For $N\geqslant1$, let
\[
\mathcal C_N(\mathbb F_q)=\mathbb F_q[t]/(t^N-1).
\]
Take $N=\pi_{k+1}$. Both $H_k$ and $H_{k+1}$ are eventually $N$-periodic
because $\pi_k\mid\pi_{k+1}$. If $B_k$ and $B_{k+1}$ are length-$N$ blocks
read from the same sufficiently late position, then
\[
p_4(t)H_{k+1}(t)=p_3(t)H_k(t)
\]
gives the following relation.

\begin{prop}\label{prop:q-block}
The neighboring periodic blocks satisfy
\[
p_4(t)B_{k+1}=p_3(t)B_k
\quad\text{in }\mathcal C_N(\mathbb F_q).
\]
The block $B_k$ is its length-$\pi_k$ block repeated $N/\pi_k$ times.
If $p_4$ is nonconstant, then its class in
$\mathcal C_N(\mathbb F_q)$ is not a unit.
\end{prop}

The last assertion follows from
\[
p_4\mid Q_{k+1}\mid t^N-1.
\]
Thus the block relation specifies the affine solution set
\[
\{B\in\mathcal C_N(\mathbb F_q):p_4B=p_3B_k\}
\]
rather than the value of an invertible circulant operator applied to $B_k$.

\subsection{Pascal's array modulo \texorpdfstring{$p$}{p}}

For
\[
g(t)=\frac{1}{1-t},
\qquad
f(t)=\frac{t}{1-t},
\]
the $k$th column has generating function
\[
H_k(t)=\frac{t^k}{(1-t)^{k+1}}.
\]
Since $\polyord((1-t)^m)=p^{\lambda_p(m)}$, its period is
\[
\pi_k=p^{\lambda_p(k+1)}.
\]
For $N=\pi_{k+1}$, sufficiently late blocks satisfy
\[
(1-t)B_{k+1}=tB_k
\quad\text{in }\mathbb F_q[t]/(t^N-1).
\]
The neighboring columns are therefore related by cyclic discrete
antidifferentiation. The failure of $1-t$ to be invertible accounts for the
nonuniqueness of a periodic antidifference when only the preceding block is
specified.


% Bibliography
\bibliographystyle{amsalpha}
\bibliography{refs}

\end{document}
